\subsection{Potenzen}
\label{subsec:fundamentals_powers}

\begin{samepage}
  \begin{DefinitionBox}{Potenz}{fundamentals_power}
    Eine Potenz besteht aus einer \emph{Basis} $a$ und einem \emph{Exponenten} $n$ und
    fasst eine wiederholte Multiplikation gleicher Faktoren zusammen:\par
    \begin{minipage}[c]{0.61\linewidth}
      \[
        a^n = \underbrace{a\cdot a\cdots a}_{n\text{-mal}},
        \qquad n\in\N,\quad n\geq1.
      \]
    Zusätzlich gilt: $a^0=1 \qquad (a\neq0)$.
    \end{minipage}\hfill
    \begin{minipage}[c]{0.3\linewidth}
      \centering
      \resizebox{\linewidth}{!}{%

        \CachedTikzFigure{content/resources/chapters/fundamentals/tikz/power_anatomy.tex}%
      }
    \end{minipage}\par\smallskip
  \end{DefinitionBox}

  \vspace{-0.5\baselineskip}
  \begin{ExampleBox}{Potenzen auswerten}{fundamentals_power_examples}
    \begin{enumerate}[label=(\alph*),itemsep=0.55em]
      \item $5^0=\Cloze{12cm}{1}$
      \item $5^3=\Cloze{12cm}{5\cdot5\cdot5=125}$
      \item $0^2=\Cloze{12cm}{0\cdot0=0}$
      \item $0^0$ ist \Cloze{12cm}{im Allgemeinen nicht definiert}.
    \end{enumerate}
  \end{ExampleBox}
\end{samepage}

\subsubsection{Potenzgesetze}
\label{subsubsec:fundamentals_power_laws}

Die drei Potenzgesetze vereinfachen Produkte, Quotienten und Potenzen von
Potenzen.

\begin{TheoremBox}{Potenzgesetze}{fundamentals_power_laws}
  \begingroup
  \setlength{\abovedisplayskip}{0.45em}
  \setlength{\abovedisplayshortskip}{0.45em}
  \setlength{\belowdisplayskip}{0pt}
  \setlength{\belowdisplayshortskip}{0pt}
  \begin{itemize}[leftmargin=*,topsep=0.35em,itemsep=0.9em]
    \item \textbf{Erstes Potenzgesetz -- bei gleicher Basis}
      \[
        a^m\cdot a^n=a^{m+n}
        \qquad\text{und}\qquad
        \frac{a^m}{a^n}=a^{m-n}\quad(a\neq0).
      \]

    \item \textbf{Zweites Potenzgesetz -- bei gleichem Exponenten}
      \[
        a^n\cdot b^n=(ab)^n
        \qquad\text{und}\qquad
        \frac{a^n}{b^n}=\left(\frac ab\right)^n\quad(b\neq0).
      \]

    \item \textbf{Drittes Potenzgesetz -- Verkettung von Potenzen}
      \[
        (a^m)^n=a^{mn}=(a^n)^m.
      \]
  \end{itemize}
  \endgroup
\end{TheoremBox}

Beim 3.Potenzgesetz ist auf richtige Klammersetzung zu achten, denn im Allgemeinen gilt:
\[
  (a^m)^n=a^{mn}\neq a^{(m^n)}=a^{m^n}.
\]

\begin{ExampleBox}{Potenzgesetze}{fundamentals_power_laws_examples}
  \begin{itemize}[leftmargin=*,itemsep=0.8em]
    \item \textbf{Erstes Potenzgesetz -- gleiche Basis}
      \begin{enumerate}[
        label=(\alph*),itemsep=0.35em,series=power_law_examples]
        \item
          \EquationBranches
            {3^4\cdot3^2}
            {\Cloze{6cm}{81\cdot9=729}}
            {(ohne Potenzgesetz)}
            {\Cloze{6cm}{3^{4+2}=3^6=729}}
            {(mit dem ersten Potenzgesetz)}
        \item
          \EquationBranches
            {\displaystyle\prod_{k=1}^{100}a^k}
            {\Cloze{6cm}{a \cdot a^2 \cdot a^3 \cdot \ldots \cdot a^{100} = \ldots}}
            {(ohne Potenzgesetz)}
            {\Cloze{6cm}{a^{\sum_{k=1}^{100}k}
              =a^{\frac{100\cdot101}{2}}=a^{5050}}}
            {(mit dem ersten Potenzgesetz)}
      \end{enumerate}

    \item \textbf{Zweites Potenzgesetz -- gleicher Exponent}
      \begin{enumerate}[
        label=(\alph*),itemsep=0.35em,resume=power_law_examples]
        \item
          \EquationBranches
            {2^3\cdot5^3}
            {\Cloze{6cm}{8\cdot125=1000}}
            {(ohne Potenzgesetz)}
            {\Cloze{6cm}{(2\cdot5)^3=10^3=1000}}
            {(mit dem zweiten Potenzgesetz)}
        \item
          \EquationBranches
            {\displaystyle\frac{6^3}{3^3}}
            {\Cloze{6cm}{\displaystyle\frac{216}{27}=8}}
            {(ohne Potenzgesetz)}
            {\Cloze{6cm}{\left(\frac63\right)^3=2^3=8}}
            {(mit dem zweiten Potenzgesetz)}
      \end{enumerate}

    \item \textbf{Drittes Potenzgesetz -- Verkettung von Potenzen}
      \begin{enumerate}[
        label=(\alph*),itemsep=0.35em,resume=power_law_examples]
        \item
          \EquationBranches
            {(2^3)^2}
            {\Cloze{6cm}{8^2=64}}
            {(ohne Potenzgesetz)}
            {\Cloze{6cm}{2^{3\cdot2}=2^6=64}}
            {(mit dem dritten Potenzgesetz)}
        \item $5^{(2^5)}=\Cloze{12cm}{5^{32}
          \qquad
          \text{\textcolor{red}{(aber $5^{(2^5)} \neq 5^{2\cdot 5}$,
                                 da 3. Potenzgesetz nicht anwendbar) }}}$
      \end{enumerate}
  \end{itemize}
\end{ExampleBox}

\begin{RemarkBox}
  Als Spezialfall des 1. Potenzgesetzes mit $m=-n$ ergibt sich
  \[
  \EquationMerge
    {}
    {\Cloze{6cm}{a^n\cdot\MathHighlight{a^{-n}}=a^{n+(-n)}=a^0}}
    {}
    {\Cloze{4cm}{a^n\cdot\MathHighlight{\frac{1}{a^n}}}}
    {\Cloze{1cm}{1}}
  \]
  Ein Vergleich der Terme zeit, dass
  \begin{ClozeArea}[8]{3}
    \[
      a^{-n}=\frac{1}{a^n}.
    \]
  \end{ClozeArea}
  Kehrwertbildung tauscht das Vorzeichen im Exponenten!
\end{RemarkBox}
